% ============================================================
% content/solutions.tex —— 2020 全国Ⅰ卷（理综物理部分）· 参考答案与详细解析
% 解析类型：source_embedded（转录自来源站 Word 版内嵌解析，独立复核并整理）
% ============================================================

\begin{answer}
\gksolhead{14}{单项选择题}
\gkanswer{答案 D}
设司机碰撞前的速度为 $v$，碰撞后速度减为零。无论是否使用安全气囊，司机的初、末速度都相同，
动量变化量均为
\[
  \Delta p=0-mv=-mv,
\]
其大小不变，故 B 错误。

由动量定理
\[
  F\Delta t=\Delta p
\]
可知，在动量变化量一定的条件下，安全气囊延长了司机与车体的作用时间 $\Delta t$，使平均作用力 $F$ 减小；
同时气囊增大了司机的受力面积，单位面积的受力进一步减小。故 A 错误、D 正确。

碰撞后汽车与司机最终都静止，司机的动能转化为内能和气囊的弹性势能，不可能“全部转换成汽车的动能”，
故 C 错误。故选 D。
\end{answer}

\begin{answer}
\gksolhead{15}{单项选择题}
\gkanswer{答案 B}
设地球质量为 $M$、半径为 $R$，物体质量为 $m$。由题意，火星质量为 $\dfrac{M}{10}$、半径为 $\dfrac{R}{2}$。
由万有引力定律，物体在星球表面受到的引力
\[
  F=\frac{GMm}{r^{2}}.
\]
同一物体在火星表面与在地球表面受到的引力之比
\[
  \frac{F_{\text{火}}}{F_{\text{地}}}
  =\frac{\dfrac{GMm}{10\left(\dfrac{R}{2}\right)^{2}}}{\dfrac{GMm}{R^{2}}}
  =\frac{\dfrac{1}{10}}{\dfrac{1}{4}}
  =0.4.
\]
故选 B。
\end{answer}

\begin{answer}
\gksolhead{16}{单项选择题}
\gkanswer{答案 B}
在秋千支架正下方（最低点），两根绳的拉力 $F_T$ 与重力 $mg$ 的合力提供向心力。设每根绳的拉力为 $F_T$，有
\[
  2F_T-mg=m\frac{v^{2}}{l}.
\]
代入 $m=\SI{50}{\kilogram}$、$l=\SI{10}{\metre}$、$v=\SI{8}{\metre\per\second}$、$g=\SI{10}{\metre\per\second\squared}$：
\[
  F_T=\frac{1}{2}\left(mg+m\frac{v^{2}}{l}\right)
      =\frac{1}{2}\left(50\times10+50\times\frac{64}{10}\right)\si{\newton}
      =\frac{1}{2}(500+320)\si{\newton}
      =410\si{\newton}.
\]
最接近 \SI{400}{\newton}，故选 B。
\end{answer}

\begin{answer}
\gksolhead{17}{单项选择题}
\gkanswer{答案 A}
电容器所带电荷量 $Q=CU_C$，充电电流
\[
  i=\frac{\Delta Q}{\Delta t}=C\frac{\Delta U_C}{\Delta t},
\]
电阻 $R$ 两端电压
\[
  U_R=iR=CR\frac{\Delta U_C}{\Delta t}.
\]
设图乙中电压的稳定值为 $U_m$，分段判断：

0--1 s：$U_C=0$ 且保持不变，$\dfrac{\Delta U_C}{\Delta t}=0$，故 $U_R=0$。

1--2 s：$U_C$ 由 $0$ 线性增大到 $U_m$，$\dfrac{\Delta U_C}{\Delta t}=\dfrac{U_m}{1\,\mathrm{s}}$，故
\[
  U_{R1}=\frac{CRU_m}{1\,\mathrm{s}},
\]
为正向恒定值。

2--3 s：$U_C=U_m$ 保持不变，故 $U_R=0$。

3--5 s：$U_C$ 由 $U_m$ 线性减小到 $0$，$\dfrac{\Delta U_C}{\Delta t}=-\dfrac{U_m}{2\,\mathrm{s}}$，故
\[
  U_{R2}=-\frac{CRU_m}{2\,\mathrm{s}},
\]
为负向恒定值，大小是 $U_{R1}$ 的一半。

即 $U_R$ 依次为零、正向恒定、零、负向恒定，且后两段的电压大小之比为 $2:1$。四个选项中只有 A 符合。
故选 A。
\end{answer}

\begin{answer}
\gksolhead{18}{单项选择题}
\gkanswer{答案 C}
粒子在磁场中做匀速圆周运动，洛伦兹力提供向心力：
\[
  qvB=m\frac{v^{2}}{r},\qquad r=\frac{mv}{qB},
\]
周期
\[
  T=\frac{2\pi r}{v}=\frac{2\pi m}{qB}.
\]
粒子在磁场中运动的时间
\[
  t=\frac{\theta}{2\pi}T=\frac{\theta m}{qB},
\]
其中 $\theta$ 为轨迹对应的圆心角。可见运动时间只由 $\theta$ 决定，与速率无关，$\theta$ 越大时间越长。

设半圆的半径为 $R$，则 $ac=R$，$ab=2R$。粒子自 $c$ 点垂直于 $ac$ 射入，轨迹圆心必在直线 $ac$ 上。
设轨迹半径为 $r$，按 $r$ 的大小分三种情况：

（1）$r\le \dfrac{R}{2}$：粒子从 $ca$ 段射出磁场，轨迹为半圆，$\theta=\pi$。

（2）$r\ge \dfrac{3R}{2}$：粒子从 $bd$ 段射出磁场，轨迹为半圆，$\theta=\pi$。

（3）$\dfrac{R}{2}<r<\dfrac{3R}{2}$：粒子从半圆弧 $ab$ 射出。$r$ 由 $\dfrac{R}{2}$ 增大时，出射点由 $a$ 沿圆弧向 $b$ 移动，
轨迹圆心角由 $\pi$ 逐渐增大；当 $r=R$ 时，轨迹圆心与 $a$ 重合，轨迹圆与边界圆半径相等、圆心距为 $R$，
两交点与两圆心构成等边三角形，出射点处的轨迹半径与 $ab$ 方向成 $60^{\circ}$ 角，故
\[
  \theta_{\max}=\pi+\frac{\pi}{3}=\frac{4\pi}{3}.
\]
$r$ 继续增大时，圆心角又逐渐减小，至 $r=\dfrac{3R}{2}$ 时回到 $\pi$。

\begin{center}
  \small 轨迹半径 $r=R$ 时的最大圆心角如图所示。
\end{center}
\gkfig{figures/s18-trajectory}

因此最大圆心角为 $\dfrac{4\pi}{3}$，最长运动时间
\[
  t_{\max}=\frac{\theta_{\max}}{2\pi}T
         =\frac{\frac{4\pi}{3}}{2\pi}\cdot\frac{2\pi m}{qB}
         =\frac{4\pi m}{3qB}.
\]
故选 C。
\end{answer}

\begin{answer}
\gksolhead{19}{多项选择题}
\gkanswer{答案 BD}
核反应遵循质量数守恒和电荷数守恒。设 $\mathrm{X}$ 的质量数为 $A$、电荷数为 $Z$，逐项检验：

A．${}_{1}^{2}\mathrm{H}+{}_{1}^{2}\mathrm{H}\to{}_{0}^{1}\mathrm{n}+\mathrm{X}_1$：
$2+2=1+A$，得 $A=3$；$1+1=0+Z$，得 $Z=2$。$\mathrm{X}_1$ 为 ${}_{2}^{3}\mathrm{He}$，不是 $\alpha$ 粒子。

B．${}_{1}^{2}\mathrm{H}+{}_{1}^{3}\mathrm{H}\to{}_{0}^{1}\mathrm{n}+\mathrm{X}_2$：
$2+3=1+A$，得 $A=4$；$1+1=0+Z$，得 $Z=2$。$\mathrm{X}_2$ 为 ${}_{2}^{4}\mathrm{He}$，即 $\alpha$ 粒子。

C．${}_{92}^{235}\mathrm{U}+{}_{0}^{1}\mathrm{n}\to{}_{56}^{144}\mathrm{Ba}+{}_{36}^{89}\mathrm{Kr}+3\mathrm{X}_3$：
$235+1=144+89+3A$，得 $A=1$；$92+0=56+36+3Z$，得 $Z=0$。$\mathrm{X}_3$ 为中子，不是 $\alpha$ 粒子。

D．${}_{0}^{1}\mathrm{n}+{}_{3}^{6}\mathrm{Li}\to{}_{1}^{3}\mathrm{H}+\mathrm{X}_4$：
$1+6=3+A$，得 $A=4$；$0+3=1+Z$，得 $Z=2$。$\mathrm{X}_4$ 为 ${}_{2}^{4}\mathrm{He}$，即 $\alpha$ 粒子。

故代表 $\alpha$ 粒子的是 B、D，选 BD。
\end{answer}

\begin{answer}
\gksolhead{20}{多项选择题}
\gkanswer{答案 AB}
由图像可知，物块在斜面顶端（$s=0$）时重力势能 $E_p=30\,\mathrm{J}$。由 $E_p=mgh$ 得
\[
  m=\frac{E_p}{gh}=\frac{30}{10\times3.0}\,\si{\kilogram}=1.0\,\si{\kilogram}.
\]
滑到斜面底端（$s=5.0\,\mathrm{m}$）时动能 $E_k=10\,\mathrm{J}$。初态机械能为 $30\,\mathrm{J}$，末态机械能为 $10\,\mathrm{J}$，
损失 $20\,\mathrm{J}$，故机械能不守恒，A 正确。

损失的机械能等于克服摩擦力做的功：
\[
  \mu mg\cos\theta\cdot s=20\,\mathrm{J}.
\]
由几何关系 $\sin\theta=\dfrac{h}{s}=\dfrac{3.0}{5.0}=0.6$，故 $\cos\theta=0.8$。代入得
\[
  \mu=\frac{20}{1.0\times10\times0.8\times5.0}=0.5,
\]
B 正确。

沿斜面方向，由牛顿第二定律
\[
  mg\sin\theta-\mu mg\cos\theta=ma,
\]
得
\[
  a=g\sin\theta-\mu g\cos\theta
   =(10\times0.6-0.5\times10\times0.8)\,\si{\metre\per\second\squared}
   =\SI{2.0}{\metre\per\second\squared}.
\]
C 中 \SI{6.0}{\metre\per\second\squared} 是未计摩擦时的加速度，故 C 错误。

物块下滑 $s'=2.0\,\mathrm{m}$ 时损失的机械能
\[
  Q=\mu mg\cos\theta\cdot s'
   =0.5\times10\times0.8\times2.0\,\mathrm{J}
   =8\,\mathrm{J},
\]
不是 $12\,\mathrm{J}$，故 D 错误。

故选 AB。
\end{answer}

\begin{answer}
\gksolhead{21}{多项选择题}
\gkanswer{答案 BC}
设导体棒 $MN$ 的速度为 $v_1$，金属框的速度为 $v_2$。金属框的 $bc$ 边与导体棒都在切割磁感线，
回路中的总电动势为
\[
  E=Bl(v_2-v_1).
\]
设导体棒的电阻为 $R$，则回路电流
\[
  I=\frac{E}{R}=\frac{Bl(v_2-v_1)}{R}.
\]
导体棒所受安培力
\[
  F_A=BIl=\frac{B^{2}l^{2}(v_2-v_1)}{R},
\]
方向向右（阻碍两者相对运动）；金属框 $bc$ 边所受安培力大小相同、方向向左。

设导体棒、金属框的质量分别为 $m_1$、$m_2$，由牛顿第二定律
\[
  F_A=m_1a_1,\qquad F-F_A=m_2a_2.
\]
初始时 $v_1=v_2=0$，$F_A=0$，故 $a_1=0$，$a_2=\dfrac{F}{m_2}$。
随着 $v_2-v_1$ 增大，$F_A$ 增大，$a_1$ 增大而 $a_2$ 减小。当 $a_1=a_2=a$ 时，相对速度不再变化，此后：
\[
  a=\frac{F}{m_1+m_2},
\]
\[
  F_A=m_1a=\frac{m_1F}{m_1+m_2}\quad(\text{恒定}),
\]
\[
  v_2-v_1=\frac{F_A R}{B^{2}l^{2}}
          =\frac{FRm_1}{B^{2}l^{2}(m_1+m_2)}\quad(\text{恒定}).
\]
\begin{center}
  \small 金属框与导体棒的速度—时间图像如图所示（图线 2 为金属框，图线 1 为导体棒）。
\end{center}
\gkfig{figures/s21-vt}
此后金属框与导体棒以相同的加速度 $a$ 做匀加速运动，速度继续增大；导体棒到 $bc$ 边的距离也继续增大。
因此金属框的加速度大小趋于恒定值、导体棒所受安培力趋于恒定值，而金属框的速度和棒到 $bc$ 边的距离都不会趋于恒定值。
故选 BC。
\end{answer}

\begin{answer}
\gksolhead{22}{实验题}
\gkanswer{答案 (1) $O$、$P$；(2) Ⅰ，$50.5$；(3) $50.0$}
（1）电压表跨接在 $O$、$P$ 两点时，电压表与 $R_x$ 并联，电流表测量的是通过 $R_x$ 和电压表的总电流，
电流表示数偏大，测得电阻 $R_{\text{测}}=\dfrac{U}{I}$ 偏小；电压表跨接在 $O$、$Q$ 两点时，电压表测量的是
$R_x$ 与电流表串联后的总电压，电压表示数偏大，测得电阻偏大。

在电流相同的情况下，跨接 $O$、$Q$ 时电压表示数更大，对应图线Ⅰ；跨接 $O$、$P$ 时电压表示数更小，对应图线Ⅱ。
故图线Ⅱ是采用电压表跨接在 $O$、$P$ 两点的方案测得的。

（2）$R_x$ 为几十欧姆，远大于电流表内阻 $R_A=\SI{0.5}{\ohm}$，又远小于电压表内阻 $R_V=\SI{1}{\kilo\ohm}$。
内接法（跨接 $O$、$Q$）的相对误差约为 $\dfrac{R_A}{R_x}\approx1\%$；
外接法（跨接 $O$、$P$）的相对误差约为 $\dfrac{R_x}{R_V}\approx5\%$。
故内接法更接近真实值，即由图线Ⅰ得到的结果更接近真实值。

由图线Ⅰ取两点 $(20.0\,\mathrm{mA},\,1.01\,\mathrm{V})$ 和 $(60.0\,\mathrm{mA},\,3.03\,\mathrm{V})$，则
\[
  R_x+R_A=\frac{\Delta U}{\Delta I}
          =\frac{3.03-1.01}{(60.0-20.0)\times10^{-3}}\,\si{\ohm}
          =\frac{2.02}{0.0400}\,\si{\ohm}
          \approx50.5\,\si{\ohm}.
\]
（3）图线Ⅰ对应电流表内接，测量值为 $R_x$ 与电流表内阻之和，故
\[
  R_x=(50.5-0.5)\,\si{\ohm}=50.0\,\si{\ohm}.
\]
\end{answer}

\begin{answer}
\gksolhead{23}{实验题}
\gkanswer{答案 (1) 大约相等；(5) $m_1gt_{12}$，$m_2\left(\dfrac{d}{\Delta t_2}-\dfrac{d}{\Delta t_1}\right)$；
(6) $0.221$，$0.212$；(7) $4$}
（1）气垫导轨水平时，滑块不受沿导轨方向的合力，做匀速直线运动，遮光片通过两个光电门的遮光时间大约相等。

（5）将砝码和砝码盘所受重力视为滑块所受拉力，则拉力 $F=m_1g$，作用时间为 $t_{12}$，拉力冲量
\[
  I=m_1gt_{12}.
\]
遮光片经过光电门 $A$、$B$ 时滑块的速度分别为
\[
  v_A=\frac{d}{\Delta t_1},\qquad v_B=\frac{d}{\Delta t_2},
\]
滑块动量改变量的大小
\[
  \Delta p=m_2v_B-m_2v_A
        =m_2\left(\frac{d}{\Delta t_2}-\frac{d}{\Delta t_1}\right).
\]
（6）代入数据：
\[
  I=1.50\times10^{-2}\times9.80\times1.50\,\si{\newton\second}
   =0.2205\,\si{\newton\second}\approx0.221\,\si{\newton\second}.
\]
\[
  v_A=\frac{1.000\times10^{-2}}{3.900\times10^{-2}}\,\si{\metre\per\second}
     =0.2564\,\si{\metre\per\second},\qquad
  v_B=\frac{1.000\times10^{-2}}{1.270\times10^{-2}}\,\si{\metre\per\second}
     =0.7874\,\si{\metre\per\second}.
\]
\[
  \Delta p=0.400\times(0.7874-0.2564)\,\si{\kilogram\metre\per\second}
          =0.2124\,\si{\kilogram\metre\per\second}
          \approx0.212\,\si{\kilogram\metre\per\second}.
\]
（7）
\[
  \delta=\left|\frac{I-\Delta p}{I}\right|\times100\%
       =\frac{|0.2205-0.2124|}{0.2205}\times100\%
       \approx3.7\%\approx4\%.
\]
\end{answer}

\begin{answer}
\gksolhead{24}{解答题}
\gkanswer{答案 (1) \SI{78}{\metre\per\second}；(2) \SI{2.0}{\metre\per\second\squared}，\SI{39}{\second}}
（1）飞机起飞离地时，升力与重力相等。设装载货物前后的质量分别为 $m_1$、$m_2$，起飞离地速度分别为 $v_1$、$v_2$，
则
\[
  m_1g=kv_1^{2},\qquad m_2g=kv_2^{2}.
\]
两式相比得
\[
  \frac{v_2^{2}}{v_1^{2}}=\frac{m_2}{m_1},
\]
故
\[
  v_2=v_1\sqrt{\frac{m_2}{m_1}}
     =66\times\sqrt{\frac{1.69\times10^{5}}{1.21\times10^{5}}}\,\si{\metre\per\second}
     \approx66\times1.182\,\si{\metre\per\second}
     \approx78\,\si{\metre\per\second}.
\]
（2）飞机从静止开始匀加速滑行，设加速度大小为 $a$、所用时间为 $t$。由匀变速直线运动公式
\[
  v_2^{2}=2as,\qquad v_2=at.
\]
代入 $v_2=78\,\si{\metre\per\second}$、$s=1521\,\si{\metre}$：
\[
  a=\frac{v_2^{2}}{2s}
   =\frac{78^{2}}{2\times1521}\,\si{\metre\per\second\squared}
   =\frac{6084}{3042}\,\si{\metre\per\second\squared}
   =\SI{2.0}{\metre\per\second\squared}.
\]
\[
  t=\frac{v_2}{a}=\frac{78}{2.0}\,\si{\second}=39\,\si{\second}.
\]
\end{answer}

\begin{answer}
\gksolhead{25}{解答题}
\gkanswer{答案 (1) $E=\dfrac{mv_0^{2}}{2qR}$；(2) $v_1=\dfrac{\sqrt{2}}{4}v_0$；(3) $0$ 或 $\dfrac{\sqrt{3}}{2}v_0$}
（1）初速度为零的粒子只受电场力，沿电场方向做匀加速直线运动，最终自 $C$ 点射出，
说明电场方向由 $A$ 指向 $C$。

$AB$ 为直径，$O$ 为圆心，$OA=OC=R$；又 $\angle OAC=\theta=60^{\circ}$，故 $\triangle AOC$ 为等边三角形，
\[
  AC=R.
\]
由动能定理
\[
  qE\cdot AC=\frac{1}{2}mv_0^{2},
\]
得
\[
  E=\frac{mv_0^{2}}{2qR}.
\]

（2）$AB$ 为直径，故 $\angle ACB=90^{\circ}$，即 $AC\perp BC$。电场沿 $AC$ 方向，等势线与 $AC$ 垂直、与 $BC$ 平行。
粒子沿 $AC$ 方向的位移越大，电场力做的功 $qEy$ 越大，动能增量越大。

作与 $BC$ 平行的直线与圆相切于 $D$ 点，与 $AC$ 的延长线交于 $P$ 点，则自 $D$ 点射出的粒子沿 $AC$ 方向的位移最大。
由几何关系
\[
  \angle PAD=30^{\circ},\qquad AP=\frac{3R}{2},\qquad DP=\frac{\sqrt{3}R}{2}.
\]
\gkfig{figures/s25-geometry}
设粒子进入电场时的速度为 $v_1$，在电场中运动的时间为 $t_1$。粒子沿 $AC$ 方向做初速度为零的匀加速运动，
垂直 $AC$ 方向做匀速运动：
\[
  a=\frac{qE}{m}=\frac{v_0^{2}}{2R},\qquad
  AP=\frac{1}{2}at_1^{2},\qquad DP=v_1t_1.
\]
由 $AP=\dfrac{1}{2}at_1^{2}$ 得
\[
  t_1=\sqrt{\frac{2AP}{a}}
     =\sqrt{\frac{3R}{v_0^{2}/(2R)}}
     =\frac{\sqrt{6}R}{v_0}.
\]
故
\[
  v_1=\frac{DP}{t_1}
     =\frac{\sqrt{3}R/2}{\sqrt{6}R/v_0}
     =\frac{\sqrt{2}}{4}v_0.
\]

（3）粒子在电场中运动时，动量变化量沿电场方向，大小为
\[
  \Delta p=qEt=mat.
\]
令 $\Delta p=mv_0$，得
\[
  at=v_0,\qquad t=\frac{v_0}{a}=\frac{2R}{v_0}.
\]
以 $A$ 为原点，粒子进入电场的方向为 $x$ 轴正方向，电场方向为 $y$ 轴正方向。粒子在 $x$ 方向做匀速运动，
在 $y$ 方向做初速度为零的匀加速运动：
\[
  x=vt,\qquad y=\frac{1}{2}at^{2}=R.
\]
圆在该坐标系中的方程为
\[
  \left(x-\frac{\sqrt{3}R}{2}\right)^{2}+\left(y-\frac{R}{2}\right)^{2}=R^{2}.
\]
代入 $y=R$：
\[
  \left(x-\frac{\sqrt{3}R}{2}\right)^{2}+\frac{R^{2}}{4}=R^{2}
  \;\Longrightarrow\;
  \left(x-\frac{\sqrt{3}R}{2}\right)^{2}=\frac{3R^{2}}{4},
\]
解得
\[
  x=0\quad\text{或}\quad x=\sqrt{3}R.
\]
由 $x=vt$：
\[
  v=0\quad\text{或}\quad v=\frac{\sqrt{3}R}{t}=\frac{\sqrt{3}}{2}v_0.
\]
即粒子进入电场时的速度应为 $0$ 或 $\dfrac{\sqrt{3}}{2}v_0$。
\end{answer}

\begin{answer}
\gksolhead{33}{选修 3-3}
\gkanswer{答案 (1) 减小；减小；小于；(2)(i) $\dfrac{2p}{3}$；(ii) $\dfrac{2}{3}$}
（1）两分子从相距无穷远到间距为 $r_1$ 的过程中，分子力表现为引力，分子力做正功，分子势能减小；
其中间距减小到 $r_2$ 的过程势能一直减小，由 $r_2$ 减小到 $r_1$ 的过程势能继续减小。
在 $r=r_1$ 处分子力为零，分子势能最小，且小于零（规定无穷远处势能为零）。
故依次填：减小；减小；小于。

（2）(i) 设乙罐中的气体等温压缩到压强为 $p$ 时体积为 $V_1$。由玻意耳定律
\[
  \frac{1}{2}p\cdot 2V=pV_1,
\]
得
\[
  V_1=V.
\]
此时两罐气体压强均为 $p$，总体积为 $V+V_1=2V$。设调配后两罐中气体的压强为 $p'$，总体积为 $V+2V=3V$，
由玻意耳定律
\[
  p\cdot 2V=p'\cdot 3V,
\]
得
\[
  p'=\frac{2}{3}p.
\]

(ii) 调配后甲罐中气体的压强为 $p'$、体积为 $V$。设这部分气体等温压缩回压强 $p$ 时体积为 $V_2$，由玻意耳定律
\[
  p'V=pV_2,
\]
得
\[
  V_2=\frac{p'}{p}V=\frac{2}{3}V.
\]
同种气体在相同温度、相同压强下，质量与体积成正比，故调配后甲罐中气体的质量与甲罐中原有气体的质量之比为
\[
  \frac{m_{\text{后}}}{m_{\text{原}}}=\frac{V_2}{V}=\frac{2}{3}.
\]
\end{answer}

\begin{answer}
\gksolhead{34}{选修 3-4}
\gkanswer{答案 (1) BCE；(2)(i) $\dfrac{l}{4}$；(ii) $\dfrac{fl}{4}$}
（1）B 中超声波被血流反射后频率发生变化，C 中列车远去时汽笛声的音调变低，E 中双星光谱随时间周期性变化，
都是波源与观察者之间存在相对运动、接收频率发生变化的现象，属于多普勒效应。
A 中“先看到闪电后听到雷声”是因为光速远大于声速；D 中声波在空气和水中传播速度不同是因为介质不同。
二者都与多普勒效应无关。故选 BCE。

（2）(i) $c$ 点到两波源 $a$、$b$ 的距离相等，波程差为零，故 $c$ 点为振动加强点。
沿 $ac$ 连线从 $c$ 向 $a$ 移动时，波程差由零逐渐增大。设距 $c$ 最近的加强点为 $d$，
$a$、$b$ 到 $d$ 的距离分别为 $r_1$、$r_2$，则该点满足
\[
  r_2-r_1=\lambda.
\]
设 $cd=s=\dfrac{3}{8}l$，则
\[
  r_1=ad=l-s=l-\frac{3}{8}l=\frac{5}{8}l.
\]
\gkfig{figures/s34-geometry}
以 $a$ 为原点、$ab$ 为 $x$ 轴，等边三角形 $abc$ 中 $c$ 的坐标为 $\left(\dfrac{l}{2},\dfrac{\sqrt{3}l}{2}\right)$，
$d$ 在 $ac$ 上且 $ad=r_1$。由几何关系
\[
  r_2^{2}=(r_1\sin60^{\circ})^{2}+(l-r_1\cos60^{\circ})^{2}.
\]
代入 $r_1=\dfrac{5}{8}l$：
\[
  r_2^{2}=\left(\frac{5l}{8}\cdot\frac{\sqrt{3}}{2}\right)^{2}
          +\left(l-\frac{5l}{8}\cdot\frac{1}{2}\right)^{2}
         =\frac{75l^{2}}{256}+\frac{121l^{2}}{256}
         =\frac{196l^{2}}{256},
\]
故
\[
  r_2=\frac{14}{16}l=\frac{7}{8}l.
\]
于是
\[
  \lambda=r_2-r_1=\frac{7}{8}l-\frac{5}{8}l=\frac{l}{4}.
\]
(ii) 由波速公式 $v=f\lambda$ 得
\[
  v=f\cdot\frac{l}{4}=\frac{fl}{4}.
\]
\end{answer}
